%% Paper 1 (PRE): Supplemental Material draft.
%% The assumptions, statements and proofs are in pre_math_body.tex (shared
%% with pre_math_report.tex); edit them there.
%% Compile: pdflatex supplement && bibtex supplement && pdflatex supplement x2
\documentclass[pre,onecolumn,superscriptaddress,amsmath,amssymb,floatfix]{revtex4-2}
\usepackage{graphicx}
\usepackage{xcolor}
\usepackage{xurl}
\usepackage{amsthm}
\usepackage[colorlinks=true,linkcolor=blue,citecolor=blue,urlcolor=blue]{hyperref}

\newtheorem{proposition}{Proposition}
\newtheorem{lemma}[proposition]{Lemma}
\newtheorem{corollary}[proposition]{Corollary}
\theoremstyle{definition}
\newtheorem{definition}[proposition]{Definition}
\newtheorem{assumption}{Assumption}
\renewcommand{\theassumption}{A\arabic{assumption}}
\theoremstyle{remark}
\newtheorem{remark}[proposition]{Remark}

\newcommand{\todo}[1]{\textcolor{red}{[TODO: #1]}}
\newcommand{\E}{\mathbb{E}}
\newcommand{\Prob}{\mathbb{P}}
\newcommand{\avg}[1]{\langle #1\rangle}
\newcommand{\figdir}{../../figures/paper/}
\newcommand{\fig}[2]{\IfFileExists{\figdir #1}{\includegraphics[width=#2]{\figdir #1}}%
  {\fbox{\parbox[c][3cm][c]{#2}{\centering\textcolor{red}{Pending figure:\\ \texttt{\detokenize{#1}}}}}}}
\renewcommand{\thesection}{S\arabic{section}}
\renewcommand{\thefigure}{S\arabic{figure}}
\renewcommand{\theequation}{S\arabic{equation}}
\renewcommand{\thetable}{S\arabic{table}}

\begin{document}
\title{Supplemental Material for ``Condensation and optimal targeting of repeated forcing in heterogeneous networks''}
\author{Jose Herrera-Diestra}
\affiliation{Department of Mathematics, Tarleton State University, Stephenville, Texas 76402, USA}
\maketitle
\tableofcontents

%%%%%%%%%%%%%%%%%%%%%%%%%%%%%%%%%%%%%%%%%%%%%%%%%%%%%%%%%%%%%%%%%%%%%%
%% S1-S7: model, assumptions, statements and proofs (shared file)
%%%%%%%%%%%%%%%%%%%%%%%%%%%%%%%%%%%%%%%%%%%%%%%%%%%%%%%%%%%%%%%%%%%%%%
%% Shared mathematical content of the PRE paper: assumptions, statements and
%% proofs. Included by pre_math_report.tex and supplement.tex (edit here only).
%% Shared math (model); included by supplement.tex and pre_math_report.tex.
%%%%%%%%%%%%%%%%%%%%%%%%%%%%%%%%%%%%%%%%%%%%%%%%%%%%%%%%%%%%%%%%%%%%%%
\section{Model and assumptions}\label{sec:model}
%%%%%%%%%%%%%%%%%%%%%%%%%%%%%%%%%%%%%%%%%%%%%%%%%%%%%%%%%%%%%%%%%%%%%%

\begin{assumption}[network]\label{as:net}
The network is a static bipartite configuration model with degree
distributions $p_k$ (class 1) and $q_k$ (class 2), with finite second moments.
\end{assumption}
\begin{assumption}[tree-like]\label{as:tree}
Clusters started by finitely many seedings are locally tree-like: partners
reached from an infected node are distinct, previously susceptible, and have
numbers of other partners drawn independently from the excess-degree
distributions $p^{(1)}_j=(j+1)p_{j+1}/\avg{k}$ (and $q^{(1)}_j$).
\end{assumption}
\begin{assumption}[transmission]\label{as:sir}
Discrete days. Each infectious node infects each susceptible partner
independently with probability $p$ per day; then every infectious node,
including those infected that day, recovers with probability $r$. Runs
continue until extinction.
\end{assumption}
\begin{assumption}[entry nodes]\label{as:entry}
Entry nodes are a subset $B$ of class 1 chosen independently of degree, with
$k\ge1$. Their importances $w_i$ are i.i.d.
\end{assumption}
\begin{assumption}[forcing]\label{as:forcing}
On day $t\le T$ each susceptible entry node $i$ is seeded with probability
$1-\exp(-\beta_{\mathrm{ext}}\alpha_tn\pi_i)$, independently. A seeding aimed at a
non-susceptible node is lost (not redirected).
\end{assumption}
\begin{assumption}[independent clusters]\label{as:indep}
Clusters started by different seedings do not overlap, and entry nodes are
not infected through the network before they are seeded.
\end{assumption}
\begin{assumption}[tail]\label{as:tail}
$w$ has a regularly varying tail, $\Prob(W\ge w)\simeq C_0w^{-(a-1)}$, $a>1$; for
results involving $\eta$ we require $0<\eta<a-1$ (finite mean response). In the degree
model ($w=k$, entry-node degrees following the class-1 law) we also require
$a>3$, so that $\kappa_1$ and hence $A$ are finite.
\end{assumption}

\begin{remark}
\ref{as:tree} fails near the epidemic threshold and for small, densely shared
cores. The error of \ref{as:indep} is of the order of the attack fraction
$C/(N_1+N_2)$; it is largest when seeding is intense and transmissibility high.
The limit $n\to\infty$ at fixed $x$ keeps this fraction finite, so the asymptotic
results are statements about the allocation problem with overlaps neglected.
Both limitations are quantified numerically (Section~\ref{sec:validation}).
\end{remark}


%% Shared math (response); included by supplement.tex and pre_math_report.tex.
%%%%%%%%%%%%%%%%%%%%%%%%%%%%%%%%%%%%%%%%%%%%%%%%%%%%%%%%%%%%%%%%%%%%%%
\section{Response to one seeding event}\label{sec:response}
%%%%%%%%%%%%%%%%%%%%%%%%%%%%%%%%%%%%%%%%%%%%%%%%%%%%%%%%%%%%%%%%%%%%%%

\begin{lemma}[transmission days and transmissibility; uses \ref{as:sir}]
$\Prob(D=d)=r(1-r)^d$, $d\ge0$, and
\[
\tau=\E[1-(1-p)^D]=\frac{(1-r)p}{p+r-pr}\le1-r.
\]
\end{lemma}
\begin{proof}
A node becomes infectious on some day and faces recovery that same day, so
$\Prob(D=0)=r$; each further day it transmits once and then recovers with
probability $r$, so $\Prob(D=d)=(1-r)^dr$. Then
$1-\tau=\E[(1-p)^D]=\sum_dr[(1-r)(1-p)]^d=r/(1-(1-r)(1-p))=r/(p+r-pr)$, and
$\tau=1-r/(p+r-pr)=(1-r)p/(p+r-pr)$. As $p\to1$, $\tau\to1-r$.
\end{proof}

\begin{lemma}[no secondary case; uses \ref{as:sir}]
For a seeded node with $k$ susceptible partners,
$\Prob(\text{no secondary case})=\E[(1-p)^{kD}]=r/(1-(1-r)(1-p)^k)$.
\end{lemma}
\begin{proof}
$\sum_dr(1-r)^d(1-p)^{kd}=r/(1-(1-r)(1-p)^k)$. Note it differs from $(1-\tau)^k$
because the $k$ edges share the same $D$.
\end{proof}

\begin{proposition}[linear response; uses \ref{as:net}--\ref{as:sir}]\label{prop:linear}
If $R^2=\tau^2\kappa_1\kappa_2<1$, a seeding into a class-1 node of degree $k$ causes
on average $s(k)=Ak$ secondary cases, $A=\tau(1+\tau\kappa_2)/(1-R^2)$.
\end{proposition}
\begin{proof}
Let $W_2$ ($W_1$) be the mean cluster size, including the root, started by a
class-2 (class-1) node infected along an edge. By \ref{as:tree} it has $\kappa_2$
($\kappa_1$) further partners on average, each infected with probability $\tau$; by
linearity of expectation (correlations through $D$ do not affect means)
$W_2=1+\tau\kappa_2W_1$, $W_1=1+\tau\kappa_1W_2$. Hence $W_1=(1+\tau\kappa_1)/(1-R^2)$,
$W_2=(1+\tau\kappa_2)/(1-R^2)$, finite iff $R^2<1$. The seeded node infects $\tau k$
partners on average, each starting a cluster of mean $W_2$: $s(k)=\tau kW_2$.
\end{proof}

\begin{remark}[general response]
In general we write the mean response of node $i$ as $Aw_i^\eta$. The degree
model is $w=k$, $\eta=1$. Other values of $\eta$ represent nonlinear responses,
for example when importance and spreading capacity are different but
correlated attributes.
\end{remark}


%% Shared math (gain); included by supplement.tex and pre_math_report.tex.
%%%%%%%%%%%%%%%%%%%%%%%%%%%%%%%%%%%%%%%%%%%%%%%%%%%%%%%%%%%%%%%%%%%%%%
\section{Targeting gain per seeding event}\label{sec:gain}
%%%%%%%%%%%%%%%%%%%%%%%%%%%%%%%%%%%%%%%%%%%%%%%%%%%%%%%%%%%%%%%%%%%%%%

\begin{definition}
$\mu_\eta(\gamma)=\sum_{i\in B}\pi_i(\gamma)w_i^\eta$. Mean secondary cases per seeding:
$A\mu_\eta(\gamma)$; mean total per seeding: $1+A\mu_\eta(\gamma)$.
\end{definition}

\begin{proposition}\label{prop:gain}
(i) $G_{\mathrm{sec}}=\mu_\eta(\gamma)/\mu_\eta(0)$ does not depend on $A$ (hence not on
$\tau$, $\kappa_1$, $\kappa_2$). (ii) $d\mu_\eta/d\gamma=\mathrm{Cov}_\pi(\ell,w^\eta)\ge0$.
(iii) $G_{\mathrm{sec}}\to w_{\max}^\eta/\mu_\eta(0)$ as $\gamma\to\infty$.
(iv) $G_{\mathrm{tot}}=(1+A\mu_\eta(\gamma))/(1+A\mu_\eta(0))$ increases with $A$ when
$\mu_\eta(\gamma)>\mu_\eta(0)$ and tends to $G_{\mathrm{sec}}$ as $A\to\infty$ ($R\to1^-$).
\end{proposition}
\begin{proof}
(ii) $\pi_i=e^{\gamma\ell_i}/Z$ gives $d\pi_i/d\gamma=\pi_i(\ell_i-\bar\ell)$,
$\bar\ell=\sum_j\pi_j\ell_j$, so $d\mu_\eta/d\gamma=\sum_i\pi_i(\ell_i-\bar\ell)w_i^\eta=\mathrm{Cov}_\pi(\ell,w^\eta)$,
non-negative since $\ell$ and $w^\eta$ are increasing functions of $w$.
(iii) As $\gamma\to\infty$, $\pi$ concentrates on $\arg\max w$. (iv)
$\partial_AG_{\mathrm{tot}}=(\mu_\eta(\gamma)-\mu_\eta(0))/(1+A\mu_\eta(0))^2$.
\end{proof}


%% Shared math (depletion); included by supplement.tex and pre_math_report.tex.
%%%%%%%%%%%%%%%%%%%%%%%%%%%%%%%%%%%%%%%%%%%%%%%%%%%%%%%%%%%%%%%%%%%%%%
\section{Repeated forcing and depletion}\label{sec:depletion}
%%%%%%%%%%%%%%%%%%%%%%%%%%%%%%%%%%%%%%%%%%%%%%%%%%%%%%%%%%%%%%%%%%%%%%

\begin{lemma}[uses \ref{as:forcing}, \ref{as:indep}]
Entry node $i$ is ever seeded with probability $1-e^{-c\pi_i}$.
\end{lemma}
\begin{proof}
While susceptible, it escapes day $t$ with probability $e^{-\beta_{\mathrm{ext}}\alpha_tn\pi_i}$,
independently across days; the product over $t\le T$ is
$e^{-\beta_{\mathrm{ext}}n\pi_i\sum_t\alpha_t}=e^{-c\pi_i}$. Under \ref{as:indep} nothing else
removes it from the susceptible pool before it is seeded.
\end{proof}

\begin{proposition}[depletion; uses \ref{as:net}--\ref{as:indep}]
\[
I(\gamma)=\sum_{i\in B}\left(1-e^{-c\pi_i}\right),\qquad
C(\gamma)=\sum_{i\in B}\left(1-e^{-c\pi_i}\right)\left(1+Aw_i^\eta\right).
\]
\end{proposition}

\begin{proposition}[optimal targeting is positive]\label{prop:pos}
Let $w$ not be constant on $B$. (i) $I(\gamma)\le I(0)$. (ii)
\[
\frac{dC}{d\gamma}=c\sum_i\pi_ie^{-c\pi_i}(\ell_i-\bar\ell)(1+Aw_i^\eta),\qquad
\left.\frac{dC}{d\gamma}\right|_0=c\,e^{-c/n}A\,\mathrm{Cov}_B(\ell,w^\eta)>0,
\]
with $\mathrm{Cov}_B$ the unweighted covariance over $B$. (iii) $C$ is non-decreasing on
$\gamma\le0$. (iv) Hence every maximizer satisfies $\gamma^*>0$, for all $c>0$. At finite $n$, $\gamma^*$
can be infinite when $c$ is small (the burden then increases up to complete
concentration); in the large-population limit it is finite
(Section~\ref{sec:optimal}). (v) $C(\gamma)=c(1+A\mu_\eta(\gamma))+O(c^2)$ as $c\to0$, non-decreasing in $\gamma$.
\end{proposition}
\begin{proof}
(i) Jensen: $\phi(\pi)=1-e^{-c\pi}$ is concave and $\sum\pi_i=1$, so
$\sum_i\phi(\pi_i)\le n\phi(1/n)=I(0)$. (ii) Differentiate and use
$d\pi_i/d\gamma=\pi_i(\ell_i-\bar\ell)$; at $\gamma=0$, $\pi_i=1/n$ and $\sum_i(\ell_i-\bar\ell)=0$
remove the constant 1. (iii) Write $dC/d\gamma=c\,\mathrm{Cov}_\pi(\ell,h)$ with
$h_i=e^{-c\pi_i}(1+Aw_i^\eta)$. For $\gamma\le0$, $\pi_i$ is non-increasing in $w_i$, so $e^{-c\pi_i}$ and
$1+Aw_i^\eta$ are non-decreasing in $w_i$, and so is $h$. The covariance of two
non-decreasing functions of $w$ is non-negative under any measure (Chebyshev's
sum inequality), so $dC/d\gamma\ge0$ on $\gamma\le0$. (iv) By (iii), $C(\gamma)\le C(0)$ for $\gamma<0$, and by
(ii), $C(\varepsilon)>C(0)$ for small $\varepsilon>0$. (v) Expand $1-e^{-c\pi}=c\pi+O(c^2)$ and use
Proposition~\ref{prop:gain}(ii).
\end{proof}


%% Shared math (transition); included by supplement.tex and pre_math_report.tex.
%%%%%%%%%%%%%%%%%%%%%%%%%%%%%%%%%%%%%%%%%%%%%%%%%%%%%%%%%%%%%%%%%%%%%%
\section{The targeting transition}\label{sec:transition}
%%%%%%%%%%%%%%%%%%%%%%%%%%%%%%%%%%%%%%%%%%%%%%%%%%%%%%%%%%%%%%%%%%%%%%
$\pi_i=e^{\gamma\ell_i}/Z$ is a Gibbs measure with inverse temperature $\gamma$ and
energies $-\ell_i$ \cite{derrida1981random}. Throughout, \ref{as:entry} and
\ref{as:tail} hold and $n\to\infty$.

\begin{lemma}[tail index]
$w^s$ ($s>0$) has tail index $(a-1)/s$: $\Prob(W^s>y)\simeq C_0y^{-(a-1)/s}$. Hence
$\E[w^s]<\infty$ iff $s<a-1$, and $w_{\max}=\max_{i\le n}w_i\sim(C_0n)^{1/(a-1)}$.
\end{lemma}

\begin{proposition}[critical exponents]
$\gamma_1^{(\eta)}=a-1-\eta$ and $\gamma_2=a-1$, with
\[
\theta_\eta(\gamma)=\begin{cases}0,&\gamma<a-1-\eta,\\ \dfrac{\gamma-(a-1-\eta)}{a-1},&a-1-\eta<\gamma<a-1,\\[4pt] \dfrac{\eta}{a-1},&\gamma>a-1.\end{cases}
\]
\end{proposition}
\begin{proof}
$\mu_\eta=\overline{w^{\gamma+\eta}}/\overline{w^\gamma}$. If both exponents are below $a-1$, both
averages converge by the strong law of large numbers, whose converse
states that the empirical mean of nonnegative i.i.d.\ variables with infinite
expectation grows without bound \cite{durrett2019probability}. Hence $\theta=0$. If $\gamma+\eta>a-1>\gamma$, the denominator
converges and the numerator, a sum of i.i.d.\ variables with tail index
$(a-1)/(\gamma+\eta)<1$, is of the order of its maximum $w_{\max}^{\gamma+\eta}$, divided by
$n$: $\mu_\eta\sim n^{(\gamma+\eta)/(a-1)-1}$. If $\gamma>a-1$, both sums are of the order of
their maxima: $\mu_\eta\sim w_{\max}^{\gamma+\eta}/w_{\max}^{\gamma}=w_{\max}^\eta\sim n^{\eta/(a-1)}$.
\end{proof}

\begin{proposition}[condensation]
For $\gamma>\gamma_2$, the ranked targeting probabilities converge in law to
Poisson--Dirichlet $\mathrm{PD}(\beta_{\mathrm t},0)$ with $\beta_{\mathrm t}=(a-1)/\gamma<1$
\cite{pitman1997two}, and $\E[Y_2]\to1-(a-1)/\gamma$. For $\gamma<\gamma_2$, $Y_2\to0$.
\end{proposition}
\begin{proof}[Sketch]
Normalized i.i.d.\ positive weights with tail index $\beta_{\mathrm t}\in(0,1)$ converge to
$\mathrm{PD}(\beta_{\mathrm t},0)$; for this law $\E[\sum_iP_i^2]=1-\beta_{\mathrm t}$. For
$\beta_{\mathrm t}>1$ (finite mean), $Y_2\le\max_i\pi_i\to0$.
\end{proof}

\begin{remark}[finite-size scaling]
For finite $n$ the transition is rounded over a window of width $O(1/\log n)$
around $\gamma_2$; we test the scaling variable $(\gamma-\gamma_2)\log n$. In the condensed phase
$Y_2$ does not self-average.
\end{remark}

\begin{proposition}[cut-off crossover]
If $w\le K_{\max}$, all moments are finite and there is no transition as $n\to\infty$.
For finite $n$ the cut-off is irrelevant while $w_{\max}(n)\ll K_{\max}$, i.e.\ for
$n\ll n_c=K_{\max}^{a-1}/C_0=(K_{\max}/k_0)^{a-1}/P_{\mathrm{tail}}$, and dominates for $n\gg n_c$.
Capped/uncapped ratios of $\mu$ and $Y_2$ are therefore functions of $n/n_c$.
\end{proposition}
\begin{proof}
$w_{\max}$ solves $n\Prob(W\ge w_{\max})\approx1$, i.e.\ $nC_0w_{\max}^{-(a-1)}\approx1$. Setting
$w_{\max}=K_{\max}$ gives $n_c$.
\end{proof}


%% Shared math (optimal); included by supplement.tex and pre_math_report.tex.
%%%%%%%%%%%%%%%%%%%%%%%%%%%%%%%%%%%%%%%%%%%%%%%%%%%%%%%%%%%%%%%%%%%%%%
\section{Optimal targeting in the large-population limit}\label{sec:optimal}
%%%%%%%%%%%%%%%%%%%%%%%%%%%%%%%%%%%%%%%%%%%%%%%%%%%%%%%%%%%%%%%%%%%%%%

\begin{lemma}[limit burden; uses \ref{as:entry}, \ref{as:tail}]
At fixed $x=c/n$ and $\gamma<\gamma_2$,
\[
\frac{C(\gamma)}n\to\mathcal C(\gamma;x)=\E\!\left[\left(1-e^{-xw^\gamma/M(\gamma)}\right)\left(1+Aw^\eta\right)\right],
\qquad M(\gamma)=\E[w^\gamma].
\]
For $\gamma>\gamma_2$, $C/n\to0$. Hence $\gamma^*_\infty(x)\le\gamma_2$.
\end{lemma}
\begin{proof}
$c\pi_i=x\,w_i^\gamma/\overline{w^\gamma}$ and $\overline{w^\gamma}\to M(\gamma)<\infty$ (LLN), so
$C/n=\overline{(1-e^{-c\pi})(1+Aw^\eta)}\to\mathcal C$. For $\gamma>\gamma_2$ a finite share of the $xn$
seedings goes to $O(1)$ nodes, each seeded at most once, while the remaining
weights are $o(1)$ per node: the number of realized seedings is $o(n)$.
\end{proof}

\begin{proposition}[large-$x$ asymptote]
As $x\to\infty$, $\gamma^*\simeq K/x$, where $K>0$ is the unique root of
$g(K)=\E\!\left[(\ell-\avg{\ell})(1+Aw^\eta)\,w^{-K}\right]=0$.
\end{proposition}
\begin{proof}
First-order condition: $\sum_ie^{-xu_i}u_i(\ell_i-\bar\ell_\gamma)(1+Aw_i^\eta)=0$ with
$u_i=n\pi_i$. Put $\gamma=K/x$; then $u_i=1+\gamma d_i+O(\gamma^2)$, $d_i=\ell_i-\avg\ell$,
$\bar\ell_\gamma\to\avg\ell$, and $xu_i=x+Kd_i+O(K\gamma)$. Dividing by $e^{-x}$ and letting
$x\to\infty$: $\sum_ie^{-Kd_i}d_i(1+Aw_i^\eta)=0$, which is $g(K)=0$ up to the positive
factor $e^{K\avg\ell}$. Existence: $g(0)=A\,\mathrm{Cov}(\ell,w^\eta)>0$, and as $K\to\infty$ the
weight $w^{-K}$ concentrates on the smallest value $w_{\min}$, where $d<0$, so $g$
becomes negative. Uniqueness: with $d=\ell-\avg\ell$,
$g'(K)=-\E[d\,\ell\,(1+Aw^\eta)w^{-K}]=-\E[d^2(1+Aw^\eta)w^{-K}]-\avg\ell\,g(K)$, so at any root
$g'(K)=-\E[d^2(1+Aw^\eta)w^{-K}]<0$. Every root is a strict down-crossing of a
continuous function, hence there is only one.
\end{proof}

\begin{proposition}[small-$x$ asymptotics]\label{prop:smallx}
As $x\to0$, $\gamma^*_\infty(x)$ is asymptotically the interior local maximum, closest to
$\gamma_2$, of
\[
G(\gamma)=-\nu(\gamma)\left[\log\tfrac1x+\log M(\gamma)\right]+\log\Gamma\big(1-\nu(\gamma)\big),
\qquad \nu(\gamma)=\frac{a-1-\eta}{\gamma},
\]
which does not depend on $A$ or $C_0$. To leading order,
\[
\gamma_2-\gamma^*_\infty(x)\simeq\frac{a-1}{\log(1/x)}.
\]
\end{proposition}
\begin{proof}
Step 1. $\mathcal C=\E[1-e^{-xw^\gamma/M}]+A\,\E[(1-e^{-xw^\gamma/M})w^\eta]$. Since
$\E[xw^\gamma/M]=x$, the first term is $x-o(x)$ for every $\gamma$ and does not move the
maximum at leading order.\\
Step 2. In the second term, $\gamma+\eta>a-1$ on the interval considered, so the
expectation is dominated by $w$ near $w_s=(M/x)^{1/\gamma}\to\infty$, where the tail
density applies: $\E[(1-e^{-xw^\gamma/M})w^\eta]\simeq C_0(a-1)\int_0^\infty w^{\eta-a}(1-e^{-xw^\gamma/M})\,dw$.\\
Step 3. Substitute $u=xw^\gamma/M$: $w=(Mu/x)^{1/\gamma}$, $dw=\gamma^{-1}(M/x)^{1/\gamma}u^{1/\gamma-1}du$,
$w^{\eta-a}dw=\gamma^{-1}(M/x)^{(\eta-a+1)/\gamma}u^{-\nu-1}du$. The integral becomes
$\gamma^{-1}(x/M)^{\nu}\int_0^\infty u^{-\nu-1}(1-e^{-u})du$, convergent for $0<\nu<1$.\\
Step 4. Integration by parts: $\int_0^\infty u^{-\nu-1}(1-e^{-u})du=\nu^{-1}\Gamma(1-\nu)$. Since
$\gamma\nu=a-1-\eta$, $\mathcal C-x\simeq\frac{AC_0(a-1)}{a-1-\eta}(x/M)^\nu\Gamma(1-\nu)$; the prefactor
does not depend on $\gamma$, so maximizing $\mathcal C$ is maximizing $G=\log[(x/M)^\nu\Gamma(1-\nu)]$.\\
Step 5. Near $\gamma_2$, $M=m_0/\delta+O(1)$ with $\delta=\gamma_2-\gamma$, $m_0=C_0(a-1)$. Using
$d\nu/d\gamma=-\nu/\gamma$ and $d\log M/d\gamma\simeq1/\delta$, the stationarity condition $G'(\gamma)=0$
reads $(\nu/\gamma)[\log(1/x)+O(\log(1/\delta))]=\nu/\delta$, so $\delta\simeq\gamma/\log(1/x)\to\gamma_2/\log(1/x)$.
The terms of relative order $\log(1/\delta)/\log(1/x)$ are not all captured by $G$ (the
neglected $\gamma$-dependence of Step 1 and the regular part of $M$ enter at the same
order), so only the leading term is claimed.
\end{proof}

\begin{remark}[assumptions of Proposition~\ref{prop:smallx}]
(a) $0<\eta<a-1$; (b) $x\to0$ such that $w_s\to\infty$; (c) $\nu$ bounded away from 1,
i.e.\ $\gamma$ away from $\gamma_1^{(\eta)}$: as $\nu\to1$, $\log\Gamma(1-\nu)\to+\infty$ because the
integral of Step 3 is then dominated by small $w$, where the tail form does
not hold, so $G$ has a spurious divergence at $\gamma_1^{(\eta)}$ and the relevant
maximum is the interior one near $\gamma_2$; (d) Step 1 treats the $\gamma$-dependence of
the importation term as subleading.
\emph{Numerical check} (Section~\ref{sec:validation}): the leading order
matches the exact limit to within about 0.07 for $\eta\ge1$ and $x\le10^{-4}$ (e.g.\
$a=3.31$, $\eta=1$: 2.184 exact against 2.185 at $x=10^{-8}$), and converges more slowly
for $\eta=0.5$, where $\nu$ is close to 1 even near $\gamma_2$. The interior maximum of $G$
exists only below an $(a,\eta)$-dependent value of $x$; where it exists it
reproduces the exact limit to within $10^{-3}$--$10^{-6}$, so $G$ also captures the
subleading corrections numerically.
\end{remark}

\paragraph{Phase diagram of $\gamma^*$.} $0<\gamma^*_\infty(x)\le\gamma_2$ for all $x$ (positivity:
Proposition~\ref{prop:pos} applied to $\mathcal C$; the bound: the Lemma above; strict
inequality observed numerically; this bound holds only as $n\to\infty$); $\gamma^*_\infty\simeq K/x$ as
$x\to\infty$; $\gamma_2-\gamma^*_\infty\simeq(a-1)/\log(1/x)$ as $x\to0$. Finite $n$ approaches the limit from
above and only logarithmically in $n$ (numerics). For the survey law at $x=10^{-4}$,
$\gamma^*=5.95$ ($n=10^4$) and 2.45 ($n=10^6$), against $\gamma^*_\infty=2.03$ and $\gamma_2=2.31$; at $x=1$,
$\gamma^*=0.55$ already for $n=100$, against 0.52. The finite-$n$ optimum stays within 0.1 of
the limit only for $x\gtrsim3/\sqrt n$ over $n=10^2$--$10^6$.


%% Shared math (threshold); included by supplement.tex and pre_math_report.tex.
%%%%%%%%%%%%%%%%%%%%%%%%%%%%%%%%%%%%%%%%%%%%%%%%%%%%%%%%%%%%%%%%%%%%%%
\section{Above threshold}\label{sec:threshold}
%%%%%%%%%%%%%%%%%%%%%%%%%%%%%%%%%%%%%%%%%%%%%%%%%%%%%%%%%%%%%%%%%%%%%%
With $h(j,y)=\E_D[(1-\tau_D(1-y))^j]$, $\tau_D=1-(1-p)^D$, and $(u_1,u_2)$ the smallest
solution of $u_1=\sum_jp^{(1)}_jh(j,u_2)$, $u_2=\sum_jq^{(1)}_jh(j,u_1)$, a seeding into a class-1
node of degree $k$ starts a major outbreak with probability $1-h(k,u_2)$, and
$\Prob(\ge1\text{ major outbreak})=1-\prod_{i\in B}[e^{-c\pi_i}+(1-e^{-c\pi_i})h(k_i,u_2)]$. This
probability also peaks at intermediate $\gamma$ (numerics).




%%%%%%%%%%%%%%%%%%%%%%%%%%%%%%%%%%%%%%%%%%%%%%%%%%%%%%%%%%%%%%%%%%%%%%
\section{Validation of the epidemic results on networks}\label{sec:validation}
%%%%%%%%%%%%%%%%%%%%%%%%%%%%%%%%%%%%%%%%%%%%%%%%%%%%%%%%%%%%%%%%%%%%%%
All comparisons are made over ensembles of independently built networks,
since the analytical results describe the configuration-model ensemble.
Standard errors are computed between networks.

\begin{figure}[h]
\fig{fig1_secondary_cases.pdf}{0.9\textwidth}
\caption{Linear response on survey-calibrated networks for ages 15--44 and
25--44 (columns). (a) Mean secondary cases against the degree of the seeded
node for five transmissibilities. Points are means over 100 networks with 30
seeding events per degree and network ($\pm2$ standard errors; degrees present in
at least five networks), and lines are $s(k)=Ak$. (b) Probability of no secondary
case, with the exact formula $r/(1-(1-r)(1-p)^k)$ as lines. The median $|z|$ is below 1
except at $R\approx0.94$, where the tree-like approximation starts to fail.}
\end{figure}

\begin{figure}[h]
\fig{fig4_major_outbreaks.pdf}{0.9\textwidth}
\caption{Above threshold, on synthetic negative-binomial networks with $\kappa\approx3$.
(a) Probability that one seeding event starts a major outbreak against the
degree of the seeded node. Points are estimates from 20 networks with 50
events each and lines are $1-h(k,u_2)$. (b) Probability of at least one major
outbreak under repeated forcing against $\gamma$ at $\tau=0.6$, with the product formula
as lines. It also peaks at an intermediate $\gamma$.}
\end{figure}

%%%%%%%%%%%%%%%%%%%%%%%%%%%%%%%%%%%%%%%%%%%%%%%%%%%%%%%%%%%%%%%%%%%%%%
\section{Finite-size scaling and asymptotics of the optimal exponent}
%%%%%%%%%%%%%%%%%%%%%%%%%%%%%%%%%%%%%%%%%%%%%%%%%%%%%%%%%%%%%%%%%%%%%%
\begin{figure}[h]
\fig{figS_transition.pdf}{\textwidth}
\caption{Targeting transition for $n=10^2$ to $10^6$ entry nodes with degrees drawn
from pure power laws ($a=2.8$, 3.31, 4), from the survey law with an uncapped
tail, and from the survey law capped at 25 (columns). (a) Participation ratio
$Y_2$ against $\gamma$, with the $n\to\infty$ limit $1-(a-1)/\gamma$ (black dashed). (b) Growth
exponent $\theta(\gamma)$ fitted over $n=10^3$ to $10^6$ (points) and the prediction (lines).
Vertical lines mark $\gamma_1=a-2$ (dotted) and $\gamma_2=a-1$ (dashed); in the capped column
they mark the values of the uncapped law for reference. Plateaus of $\theta$ at
$n=10^6$ (fit against theory): 0.548 and 0.556 ($a=2.8$), 0.422 and 0.433 ($a=3.31$),
0.339 and 0.333 ($a=4$), 0.428 and 0.433 (survey), 0.017 (capped).}
\end{figure}

\begin{figure}[h]
\fig{figS_scaling_Y2_collapse.pdf}{\textwidth}\\[4pt]
\fig{figS_scaling_mu.pdf}{\textwidth}
\caption{Finite-size scaling. Top, participation ratio $Y_2$ against $(\gamma-\gamma_2)\ln n$;
curves for $n=10^2$ to $10^6$ collapse near zero. Bottom, targeted mean degree
$\mu(\gamma)/\mu(0)$ for the same $n$, with $\gamma_1$ (dotted) and $\gamma_2$ (dashed).}
\end{figure}

\begin{figure}[h]
\fig{figS_scaling_gamma_star_limit.pdf}{0.95\textwidth}
\caption{Optimal exponent $\gamma^*$ against $x=c/n$ for $A=2$. Colored curves are
medians for $n=10^2$ to $10^6$, the black curve the $n\to\infty$ limit, the dashed
black curve $K/x$ and the grey dashed line $\gamma_2$. Finite $n$ approaches the limit
from above and only logarithmically in $n$.}
\end{figure}

\begin{figure}[h]
\fig{figS_gen_gamma_star_eta.pdf}{0.95\textwidth}
\caption{Exact $n\to\infty$ limit of $\gamma^*$ (black) against the small-$x$ formula
(maximum of $G$, red dashed, shown where its interior maximum exists), the
leading order $\gamma_2-(a-1)/\ln(1/x)$ (dotted, $x\le10^{-2}$) and $K/x$ (blue dashed,
$x\ge0.3$), for response exponents $\eta=0.5$, 1, 2 (rows) and four tail laws
(columns), with $A=2$. Horizontal dotted lines mark $\gamma_2=a-1$. The panel $a=2.8$,
$\eta=2$ is empty because $\eta\ge a-1$.}
\end{figure}

\begin{table}[h]
\caption{Large-$x$ constant $K$ ($A=2$) and the small-$x$ regime: exact limit
$\gamma^*_\infty$, leading order $\gamma_2-(a-1)/\ln(1/x)$, and the maximum of $G$ (``--'': no
interior maximum, i.e.\ asymptotic regime not yet reached).}
\begin{ruledtabular}
\begin{tabular}{lccc ccc ccc}
 & & & & \multicolumn{3}{c}{$x=10^{-8}$} & \multicolumn{3}{c}{$x=10^{-4}$} \\
Law & $a$ & $\eta$ & $K$ & exact & leading & $G$ & exact & leading & $G$ \\
\hline
power & 2.80 & 0.5 & 0.380 & 1.694 & 1.702 & 1.693 & 1.558 & 1.605 & -- \\
power & 2.80 & 1 & 0.822 & 1.708 & 1.702 & 1.708 & 1.618 & 1.605 & 1.616 \\
power & 3.31 & 0.5 & 0.372 & 2.156 & 2.185 & 2.145 & 1.950 & 2.059 & -- \\
power & 3.31 & 1 & 0.799 & 2.184 & 2.185 & 2.184 & 2.048 & 2.059 & 2.036 \\
power & 3.31 & 2 & 1.734 & 2.195 & 2.185 & 2.195 & 2.092 & 2.059 & 2.091 \\
power & 4.00 & 0.5 & 0.365 & 2.768 & 2.837 & -- & 2.453 & 2.674 & -- \\
power & 4.00 & 1 & 0.779 & 2.820 & 2.837 & 2.818 & 2.607 & 2.674 & -- \\
power & 4.00 & 2 & 1.691 & 2.841 & 2.837 & 2.840 & 2.686 & 2.674 & 2.683 \\
survey & 3.31 & 0.5 & 0.375 & 2.152 & 2.185 & 2.138 & 1.936 & 2.059 & -- \\
survey & 3.31 & 1 & 0.810 & 2.181 & 2.185 & 2.181 & 2.034 & 2.059 & 2.018 \\
survey & 3.31 & 2 & 1.758 & 2.192 & 2.185 & 2.192 & 2.082 & 2.059 & 2.081 \\
\end{tabular}
\end{ruledtabular}
\end{table}

%%%%%%%%%%%%%%%%%%%%%%%%%%%%%%%%%%%%%%%%%%%%%%%%%%%%%%%%%%%%%%%%%%%%%%
\section{Survey-calibrated networks used in the illustration}
%%%%%%%%%%%%%%%%%%%%%%%%%%%%%%%%%%%%%%%%%%%%%%%%%%%%%%%%%%%%%%%%%%%%%%
\paragraph{Degrees.} Degrees are numbers of opposite-sex partners in the past
twelve months from the 2006--2008 National Survey of Family Growth
\cite{chandra2011sexual} (Table~\ref{tab:nsfg}); the class ``0'' combines never
and not in the past year, and non-response is removed. The class of four or more
partners is spread over $k=4,\dots,25$ with $P(k)\propto k^{-(\alpha+1)}$, where $\alpha$ is the
exponent of the cumulative distribution of past-year partners
\cite{liljeros2001web}, 2.31 for men and 2.54 for women.

\begin{table}[h]
\caption{Past-year opposite-sex partners (\%), NSFG 2006--2008.}
\label{tab:nsfg}
\begin{ruledtabular}
\begin{tabular}{lccccc}
 & 0 & 1 & 2 & 3 & 4+ \\ \hline
Men 25--44 & 8.7 & 76.1 & 7.1 & 2.9 & 5.2 \\
Women 25--44 & 8.2 & 82.7 & 5.4 & 1.7 & 1.9 \\
Men 15--44 & 18.0 & 63.3 & 8.7 & 3.9 & 6.1 \\
Women 15--44 & 17.3 & 69.6 & 7.7 & 2.5 & 2.9 \\
\end{tabular}
\end{ruledtabular}
\end{table}

\paragraph{Balancing.} Men report more partners than women, a discrepancy
attributed to reporting differences and to under-sampling of very active
people \cite{brewer2000prostitution,mitchell2019why}. In the household
network the excess is removed from men with at least two partners. A man of
degree $k$ loses at most $k-1$ stubs, chosen in proportion to $k-1$. Women's
distribution and men's share without partners are then reproduced exactly;
men's share with one partner rises (76\% to 80\% for ages 25--44). Repeated
pairings of the same couple are merged. Across 100 realizations (ages
25--44), $\kappa_1\in[0.51,0.82]$, $\kappa_2\in[0.54,1.07]$ and $\sqrt{\kappa_1\kappa_2}\le0.92$.

\paragraph{Entry nodes.} Each man is bisexual with probability 0.02
\cite{copen2016sexual}; entry nodes are bisexual men with at least one female
partner (about 90 per $5\times10^3$ men).

\paragraph{Commercial-sex core.} Following Brewer et al.\ \cite{brewer2000prostitution},
who showed that sex workers missing from household surveys account for the
men--women gap in reported partners, the men's surplus is assigned to sex-worker
nodes instead of being trimmed. Sex workers are $0.05\%$ of women, the measured
prevalence of full-time-equivalent sex workers, with log-normal weights
matching the reported median and mean numbers of partners. Clients are chosen among men
with several partners with weight $(k-1)^q$, $q=8$, sampled without replacement
with exponential keys \cite{efraimidis2006weighted}, which gives about 4.6\% of men as
clients. The implied number of partners per sex worker (mean about 300 per
year) agrees with the independently measured values.

\begin{figure}[h]
\fig{figS1_network_calibration.pdf}{0.9\textwidth}
\caption{Survey calibration. (a) Degree distributions of the model networks
(20 networks pooled) against the survey, by sex and age group. (b) Excess
degrees $\kappa$ across 100 networks per age group. For ages 25--44, $\kappa$ ranges from
0.51 to 1.07 and $\sqrt{\kappa_1\kappa_2}\le0.92$; for ages 15--44 it spreads around 1.}
\end{figure}

\begin{figure}[h]
\fig{fig5_core_network.pdf}{0.9\textwidth}
\caption{Illustration with a commercial-sex core ($5\times10^4$ people per sex, 10
networks). (a) Probability that one seeding event leads to an outbreak
reaching 1\% of the population, for an entry node that is a client of the core
and for other entry nodes, against $\tau$. (b) The same probability under
repeated forcing against $\gamma$, for $c=10$ and 100.}
\end{figure}

%%%%%%%%%%%%%%%%%%%%%%%%%%%%%%%%%%%%%%%%%%%%%%%%%%%%%%%%%%%%%%%%%%%%%%
\section{Numerical methods}
%%%%%%%%%%%%%%%%%%%%%%%%%%%%%%%%%%%%%%%%%%%%%%%%%%%%%%%%%%%%%%%%%%%%%%
\begin{itemize}
\item \emph{Transmission.} $\rho=1/14$ per day ($r=0.069$); for a target $\tau$,
$p=\tau r/[(1-r)(1-\tau)]$. Runs continue until no node is infectious. The forcing
profile $\alpha_t$ is a 17-week incidence curve normalized to a maximum of 1
($\sum_t\alpha_t=48.4$) \cite{nycdohmh2022mpox}.
\item \emph{Linear response.} 100 networks per degree distribution
($5\times10^3$ people per sex), 30 seeding events per degree and network,
$\tau\in\{0.2,0.4,0.6,0.8,0.9\}$.
\item \emph{Repeated forcing.} 50 networks, 10 runs per combination of
$c\in\{10,50,200\}$, $\tau\in\{0.6,0.9\}$ and $\gamma\in\{-2,-1.5,\dots,6\}$; theory evaluated per network.
\item \emph{Optimal exponent on networks.} 100 networks per distribution;
$c=10^{0},10^{0.1},\dots,10^{3.5}$; $\gamma\in[-1,30]$ in steps of 0.05; curves flat to a
relative tolerance $10^{-9}$ are reported as not identifiable.
\item \emph{Scaling.} Entry-node degrees drawn i.i.d.\ without cap for $n=10^2$
to $10^6$ (200--2000 samples per $n$); all sums over distinct degree values;
$\theta$ fitted as the slope of the mean of $\ln\mu$ against $\ln n$ for $n\ge10^3$.
\item \emph{Limit curve.} Expectations summed exactly to $k=10^6$ with the tail
integrated analytically; $1-e^{-y}$ evaluated as $-\mathrm{expm1}(-y)$; finite-$n$
optimum on a grid with step 0.005 below $\gamma=1$ and 0.01 above.
\item \emph{Major outbreaks.} A run is major once it reaches half the
theoretical major-outbreak size (minimum 100 cases).
\item \emph{Core network.} 10 networks, $5\times10^4$ people per sex; 300 single
seeding events per entry-node type and 30 forced runs per combination; a
large outbreak is one reaching 1000 cases ($\approx1\%$ of the population).
\end{itemize}
All code, simulation outputs and figure scripts are available at
\url{https://github.com/jdiestra1977/degree-targeted-importation} (archived at \todo{Zenodo DOI}).

\bibliography{../references}
\end{document}
